\begin{figure}[htbp]
\centering
\begin{tikzpicture}[
 box/.style={draw=linkblue!65,rounded corners=2pt,align=center,
 fill=linkblue!3,text width=5.1cm,inner sep=6pt,font=\small},
 edge/.style={-{Stealth[length=2mm]},draw=linkblue!80,thick}]
\node[box] (steps) at (0,0)
 {Positive traces and ordinary steps\\Exact line-bundle descent};
\node[box] (terminal) at (6.2,0)
 {Discrepancies and analytic cycle ranks\\Terminal fourfold termination\\with real boundary};
\node[box] (scaling) at (0,-2)
 {Ordinary scaling from $X$\\Vanishing thresholds\\Stabilized auxiliary extraction};
\node[box] (limit) at (6.2,-2)
 {The limiting terminal pair\\A finite small program\\to a nef ordinary adjoint};
\node[box] (gap) at (0,-4.1)
 {Rational nef decomposition\\Fixed Cartier indices\\and a positive intersection gap};
\node[box] (perturb) at (6.2,-4.1)
 {A nearby generalized program\\Crepant rational adjoints\\Descent to the original flip base};
\node[box,text width=11.3cm] (finite) at (3.1,-6.1)
 {The original negative ray becomes trivial: contradiction\\
 A finite ordinary program and its nef or Mori endpoint
 (Theorem~\ref{thm:finite})};
\draw[edge] (steps)--(scaling);
\draw[edge] (terminal)--(limit);
\draw[edge] (scaling)--(limit);
\draw[edge] (limit.south west)--(gap.north east);
\draw[edge] (gap)--(perturb);
\draw[edge] (limit)--(perturb);
\draw[edge] (perturb)--(finite);
\end{tikzpicture}
\caption{The finite-program argument. The ordinary sequence starts on
$X$. Extraction and the limiting and nearby programs are auxiliary
comparisons. Exact bundle descent preserves the indices used in the
positive intersection gap; the last descent occurs on a common
resolution, as in Figure~\ref{fig:perturbation-descent}.}
\label{fig:dependencies}
\end{figure}
